\documentclass[11pt,letterpaper]{article}
\usepackage[margin=2.5cm]{geometry}
\usepackage[utf8]{inputenc}
\usepackage[T1]{fontenc}
\usepackage{lmodern}
\usepackage{amsmath,amssymb,amsthm,mathtools,bm}
\usepackage{enumitem}
\usepackage{booktabs}
\usepackage{tikz}
\usetikzlibrary{arrows.meta,patterns,positioning,calc}
\usepackage{pgfplots}
\pgfplotsset{compat=1.16}
\usepackage{titlesec}
\usepackage[hidelinks]{hyperref}

\setlength{\parindent}{0pt}
\setlength{\parskip}{0.5em}
\titleformat{\section}{\large\bfseries}{\thesection.}{0.6em}{}
\titleformat{\subsection}{\normalsize\bfseries}{\thesubsection.}{0.6em}{}

\theoremstyle{definition}
\newtheorem{definition}{Definition}[section]
\newtheorem{example}{Example}[section]
\newtheorem{exercise}{Exercise}[section]
\newtheorem{theorem}{Theorem}[section]
\newtheorem{remark}{Remark}[section]

\newcommand{\R}{\mathbb{R}}
\newcommand{\N}{\mathbb{N}}
\newcommand{\Z}{\mathbb{Z}}
\newcommand{\U}{\mathcal{U}}
\newcommand{\E}{\mathbb{E}}
\newcommand{\Var}{\operatorname{Var}}
\newcommand{\Cov}{\operatorname{Cov}}
\newcommand{\Normal}{\mathcal{N}}
\newcommand{\dd}{\,\mathrm{d}}
\newcommand{\T}{^{\top}}
\renewcommand{\P}{P}
\newcommand{\set}[1]{\left\{#1\right\}}
\newcommand{\given}{\mid}
\newcommand{\note}[1]{\textit{#1}}

% --- Homework exercises (tinted box) and highlighted derivation steps ---
\usepackage{tcolorbox}
\tcbuselibrary{skins,breakable}
\definecolor{hwfill}{RGB}{235,243,252}
\definecolor{hwframe}{RGB}{31,78,140}
\definecolor{newterm}{RGB}{196,78,0}
\newtcolorbox{homework}[1]{enhanced, breakable,
  colback=hwfill, colframe=hwframe, coltitle=white,
  fonttitle=\bfseries, boxrule=0.6pt, arc=2pt,
  left=6pt, right=6pt, top=4pt, bottom=4pt,
  before skip=10pt, after skip=10pt,
  title={Exercise (Homework #1)}}
\newcommand{\soln}{\par\smallskip\noindent\textbf{Solution.}\ }
% \new{...} marks what is new in a derivation step: a substituted or simplified term.
\newcommand{\new}[1]{{\color{newterm}#1}}

\title{Lecture 5: Bayesian Statistics}
\author{}\date{}
\newcommand{\vy}{\vec y}
\newcommand{\vmu}{\vec\mu}
\newcommand{\one}{\mathbb{1}}

\begin{document}
\maketitle

\section{Bayes' rule}
Let us consider some observed data
\[
y = (y_1, \dots, y_n)^{\top}.
\]
Bayesian statistics has two main goals:
\begin{enumerate}[label=(\roman*)]
\item To make inference about some unknown (unobserved) quantities $\theta = (\theta_1, \theta_2, \dots)^{\top}$, which correspond to population parameters, physical quantities, etc.
\item Prediction of new or unobserved data $\tilde y$.
\end{enumerate}

\subsection{Posterior distribution $p(\theta \mid y)$}
\[
p(\theta \mid Y = y) = \frac{p(Y = y \mid \theta)\,p(\theta)}{p(Y = y)},
\]
or more compactly
\[
p(\theta \mid y) = \frac{p(y \mid \theta)\,p(\theta)}{p(y)}.
\]
This is
\[
p(\theta \mid y) = \frac{p(y \mid \theta)\,p(\theta)}{\sum_{\theta'} p(y \mid \theta')\,p(\theta')}
\]
for discrete variables, or
\[
f(\theta \mid y) = \frac{f(y \mid \theta)\,f(\theta)}{\int \mathrm{d}\theta'\,f(y \mid \theta')\,f(\theta')}
\]
for continuous variables. I will write $p(\cdot)$ for both discrete and continuous variables.

$p(\theta)$ is called the \emph{prior distribution}; it reflects what we know before seeing the data. $p(y \mid \theta)$ is the \emph{likelihood} (or data distribution); it reflects the data we expect to see for each setting of parameters.

\subsection{Covariates}
Covariates are independent variables (not in the sense of statistically independent). For the supervised case, we also have covariates $X = \set{\vec x_1, \dots, \vec x_n}$, with $\vec x_i \in \R^k$. $X$ is usually written as an $n\times k$ matrix
\[
X = \begin{pmatrix}
x_{11} & x_{12} & \cdots & x_{1k}\\
x_{21} & x_{22} & \cdots & x_{2k}\\
\vdots & & & \vdots\\
x_{n1} & x_{n2} & \cdots & x_{nk}
\end{pmatrix},
\]
where each row corresponds to a vector $x_i \in \R^k$, such that
\[
x_i = (x_{i1}\ x_{i2}\ \cdots\ x_{ik}) \qquad\text{and}\qquad x = (x_1, \dots, x_n).
\]

Covariates appear in regression models. This is, models where we have
\[
Y = f_\theta(X) + \varepsilon,
\]
with $\varepsilon$ a random error. These models, also called \emph{statistical models}, aim to capture input--output relationships. In reality, most of the time there is no actual deterministic relationship between $Y$ and $X$, because there might be other unmeasured variables (for example $Z$, so that $Y = f(X, Z)$) contributing to $Y$, including measurement errors.

Regression models assume $\E(\varepsilon \mid X = x_0, \theta = \theta_0) = 0$. Thus
\begin{align*}
\E[Y \mid x_0, \theta_0] &= \E[f_\theta(X) + \varepsilon \mid x_0, \theta_0] = \E[f_\theta(X) \mid x_0, \theta_0] + \underbrace{\E[\varepsilon \mid x_0, \theta_0]}_{0},\\
\E[f_\theta(X) \mid x_0, \theta_0] &= \E[f_{\theta_0}(x_0) \mid x_0, \theta_0] = f_{\theta_0}(x_0).
\end{align*}
In that case, we need a posterior for the parameters given the data $(y, x)$:
\[
p(\theta \mid y, x) = \frac{p(y, x \mid \theta)\,p(\theta)}{p(y, x)}.
\]
Since $p(x, \theta) = p(x \mid \theta)\,p(\theta) = p(\theta \mid x)\,p(x)$, and normally $p(x \mid \theta) = p(x)$ (equivalently $p(\theta \mid x) = p(\theta)$), we have
\begin{align*}
p(y, x) &= p(y \mid x)\,p(x),\\
p(y, x \mid \theta) &= p(y \mid x, \theta)\,p(x \mid \theta) = p(y \mid x, \theta)\,p(x).
\end{align*}
Thus
\[
p(\theta \mid y, x) = \frac{p(y \mid x, \theta)\,p(x)\,p(\theta)}{p(y \mid x)\,p(x)}
= \frac{p(y \mid x, \theta)\,p(\theta)}{p(y \mid x)}
= \frac{p(y \mid x, \theta)\,p(\theta)}{\int \mathrm{d}\theta'\,p(y \mid x, \theta')\,p(\theta')}.
\]

%=====================================================
\section{Predictive distributions}

\subsection{Prior predictive distribution}
\[
p(y) = \int \mathrm{d}\theta\,p(y, \theta) = \int \mathrm{d}\theta\,p(y \mid \theta)\,p(\theta).
\]
This is
\[
p(y_1, \dots, y_n) = \int \mathrm{d}\theta_1\int \mathrm{d}\theta_2\cdots\int \mathrm{d}\theta_k\,p(y_1, \dots, y_n \mid \theta_1, \dots, \theta_k)\,p(\theta_1, \dots, \theta_k).
\]

\subsection{Posterior predictive distribution}
Let us suppose now that we want to predict a new $\tilde Y$ (currently unobserved) outcome, e.g.\ $\tilde Y = (\tilde Y_1)$, given the previous data $y_1, \dots, y_n$. This is
\[
p(\tilde Y_1 \mid y_1, \dots, y_n) = \int \mathrm{d}\theta\,p(\tilde Y_1 \mid y_1, \dots, y_n, \theta_1, \dots, \theta_k)\,p(\theta_1, \dots, \theta_k \mid y_1, \dots, y_n),
\]
which can be written as
\[
p(\tilde Y \mid y) = \int \mathrm{d}\theta\,p(\tilde Y \mid \theta, y)\,p(\theta \mid y).
\]
More generally, we can wonder about the chances of certain values $\tilde Y = (\tilde Y_1, \tilde Y_2, \dots)$.

It is usually assumed that $\tilde Y \perp Y \mid \theta$. This is,
\[
p(y, \tilde y \mid \theta) = p(y \mid \theta)\,p(\tilde y \mid \theta).
\]
Let us consider
\[
p(y, \tilde y) = p(\tilde y \mid y)\,p(y), \qquad p(y, \tilde y \mid \theta) = p(\tilde y \mid y, \theta)\,p(y \mid \theta).
\]
Then
\[
p(y \mid \theta)\,p(\tilde y \mid \theta) = p(\tilde y \mid y, \theta)\,p(y \mid \theta)
\ \Rightarrow\
p(\tilde y \mid y, \theta) = p(\tilde y \mid \theta).
\]
Therefore
\[
p(\tilde y \mid y) = \int \mathrm{d}\theta\,p(\tilde y \mid \theta)\,p(\theta \mid y).
\]

\begin{remark}
$p(\tilde Y \mid y, \theta) = p(\tilde Y \mid \theta)$ is a modelling assumption. That is no longer the case for time series. That is, generally
\[
p(y_{n+1} \mid y, \theta) \neq p(y_{n+1} \mid \theta),
\]
where $\tilde Y_1 = Y_{n+1}$. Thus, with the shorter notation $y = (y_1, \dots, y_n)^{\top} = y_{1:n}$,
\[
p(y_{n+1} \mid y_{1:n}) = \int \mathrm{d}\theta\,p(y_{n+1} \mid y_{1:n}, \theta)\,p(\theta \mid y_{1:n}).
\]
\end{remark}

\paragraph{With covariates.}
Let us consider now that there are covariates involved and we want to find
\[
p(\tilde Y \mid \tilde x, x, y) = \int \mathrm{d}\theta\,p(\tilde Y \mid \theta, \tilde x, x, y)\,p(\theta \mid \tilde x, x, y).
\]
Usually $p(\theta)$ is updated to $p(\theta \mid x, y)$ using previously observed data, and it is assumed not to depend on $\tilde x$. Then
\[
p(\theta \mid \tilde x, x, y) = p(\theta \mid x, y).
\]
It is also assumed that
\[
p(\tilde y \mid \theta, \tilde x, x, y) = p(\tilde y \mid \theta, \tilde x)
\]
for any $\theta$ fixed. This is, once $\theta$ is known, past data $(x, y)$ gives no additional information about $\tilde y$ beyond what $\theta$ and $\tilde x$ encode. Combined with $p(\theta \mid \tilde x, x, y) = p(\theta \mid x, y)$, this gives
\[
p(\tilde Y \mid \tilde x, x, y) = \int \mathrm{d}\theta\,p(\tilde Y \mid \theta, \tilde x)\,p(\theta \mid x, y).
\]

%=====================================================
\section{Modelling}
Bayesian modelling can be framed in an ice-cream factory analogy. We can imagine that there are $k$ knobs whose positions correspond to values $\theta_1, \dots, \theta_k$. For example, $\theta_1$ can be the amount of vanilla, $\theta_2$ the amount of chocolate, $\theta_3$ the size of the container, etc. Thus, from the modelling standpoint, $p(\theta)$ represents the modeller's uncertainty about the unknown positions of the knobs before inspecting any ice cream.

Thus, the factory produces ice creams $Y = (Y_1, \dots, Y_n)^{\top}$ for some fixed knob positions. This is $Y \mid \theta$. However, because of uncontrolled production variability, ice creams can still differ. The likelihood $p(Y \mid \theta)$, also called \emph{sampling distribution}, describes this variability.

With $Y = (Y_1, \dots, Y_n)^{\top}$ and $y = (y_1, \dots, y_n)^{\top}$: after observing ice creams $y$, we can update our belief about the parameters $\theta$ as
\[
p(\theta \mid y) = \frac{p(y \mid \theta)\,p(\theta)}{p(y)}.
\]
Of course, the factory needs some raw material $x$ to produce ice creams. Thus
\[
p(\theta \mid y, x) = \frac{p(y \mid x, \theta)\,p(\theta)}{\int \mathrm{d}\theta'\,p(y \mid x, \theta')\,p(\theta')}.
\]

%=====================================================
\section{Conjugate priors}
There are pairs (prior, likelihood) for which we can compute the posterior in closed form.

A family $\mathcal{F}$ of distributions for $\theta$ is \emph{conjugate} to a likelihood $p(y \mid \theta)$ if
\[
p(\theta) \in \mathcal{F} \ \Rightarrow\ p(\theta \mid y) \in \mathcal{F}.
\]
Any prior in $\mathcal{F}$ is called a \emph{conjugate prior} to that likelihood.

%-----------------------------------------------------
\subsection{The Beta--Binomial model}
Let us consider the binomial likelihood
\[
p(y \mid \theta) = \operatorname{Bin}(y \mid n, \theta) = \binom{n}{y}\theta^{y}(1-\theta)^{n-y},
\]
where $n$ is the number of trials, $y$ the number of successes and $\theta$ the probability of a success (the likelihood gives meaning to $\theta$).

Let us also consider now a prior
\[
p(\theta) = \operatorname{Beta}(\theta \mid \alpha, \beta) = \underbrace{\frac{\Gamma(\alpha+\beta)}{\Gamma(\alpha)\Gamma(\beta)}}_{1/B(\alpha,\beta)}\theta^{\alpha-1}(1-\theta)^{\beta-1},
\qquad
\frac{1}{B(\alpha,\beta)} = \frac{1}{\int_0^1 t^{\alpha-1}(1-t)^{\beta-1}\dd t},
\]
for $\alpha > 0$ and $\beta > 0$, or $p(\theta) \propto \theta^{\alpha-1}(1-\theta)^{\beta-1}$. For example, for $\alpha = 1$ and $\beta = 1$, $p(\theta) \propto 1$. This is,
\[
p(\theta) = \operatorname{Beta}(\theta \mid 1, 1) = \operatorname{Uniform}(0,1).
\]
Notice that
\[
\E[\theta] = \frac{\alpha}{\alpha + \beta}, \qquad \Var(\theta) = \frac{\alpha\beta}{(\alpha+\beta)^2(\alpha+\beta+1)}.
\]
Then the posterior $p(\theta \mid y)$ is
\[
p(\theta \mid y) \propto \theta^{y}(1-\theta)^{n-y}\cdot\theta^{\alpha-1}(1-\theta)^{\beta-1}
= \theta^{\new{(y+\alpha)}-1}(1-\theta)^{\new{(\beta+n-y)}-1},
\]
so
\[
p(\theta \mid y) = \operatorname{Beta}(\theta \mid \new{y + \alpha}, \new{\beta + n - y}).
\]
One nice thing about this posterior is that the expectation value and the variance can be calculated straight away:
\[
\E[\theta \mid y] = \frac{y + \alpha}{\alpha + \beta + n},
\qquad
\Var[\theta \mid y] = \frac{(y+\alpha)(\beta + n - y)}{(\alpha+\beta+n)^2(\alpha+\beta+n+1)}.
\]

\begin{homework}{3}
Let $\theta \sim \operatorname{Beta}(\alpha, \beta)$, with PDF
\[
p(\theta) = \frac{1}{B(\alpha,\beta)}\theta^{\alpha-1}(1-\theta)^{\beta-1}, \qquad 0 < \theta < 1,\ \alpha > 0,\ \beta > 0,
\]
where $B(\alpha,\beta) = \int_0^1 t^{\alpha-1}(1-t)^{\beta-1}\dd t = \dfrac{\Gamma(\alpha)\Gamma(\beta)}{\Gamma(\alpha+\beta)}$. Show that
\begin{enumerate}[label=(\roman*)]
\item the distribution is normalized, $\int_0^1 p(\theta)\dd\theta = 1$;
\item $\E[\theta] = \dfrac{\alpha}{\alpha+\beta}$;
\item $\Var(\theta) = \dfrac{\alpha\beta}{(\alpha+\beta)^2(\alpha+\beta+1)}$.
\end{enumerate}

\soln
\begin{enumerate}[label=(\roman*)]
\item By the definition of $B(\alpha,\beta)$,
\[
\int_0^1 p(\theta)\dd\theta = \frac{1}{B(\alpha,\beta)}\int_0^1 \theta^{\alpha-1}(1-\theta)^{\beta-1}\dd\theta = \frac{\new{B(\alpha,\beta)}}{B(\alpha,\beta)} = 1.
\]
\item Multiplying by $\theta$ only raises the first exponent by one, so the integral is again a Beta function:
\[
\E[\theta] = \frac{1}{B(\alpha,\beta)}\int_0^1 \theta^{\new{\alpha}}(1-\theta)^{\beta-1}\dd\theta = \frac{B(\new{\alpha+1},\beta)}{B(\alpha,\beta)}
= \frac{\Gamma(\alpha+1)\,\Gamma(\alpha+\beta)}{\Gamma(\alpha)\,\Gamma(\alpha+\beta+1)} = \new{\frac{\alpha}{\alpha+\beta}},
\]
using $\Gamma(z+1) = z\,\Gamma(z)$ twice, the $\Gamma(\beta)$ factors having cancelled.
\item In the same way,
\[
\E[\theta^2] = \frac{B(\new{\alpha+2},\beta)}{B(\alpha,\beta)} = \frac{\Gamma(\alpha+2)\,\Gamma(\alpha+\beta)}{\Gamma(\alpha)\,\Gamma(\alpha+\beta+2)} = \new{\frac{\alpha(\alpha+1)}{(\alpha+\beta)(\alpha+\beta+1)}},
\]
so
\begin{align*}
\Var(\theta) &= \E[\theta^2] - \E[\theta]^2 = \frac{\alpha(\alpha+1)}{(\alpha+\beta)(\alpha+\beta+1)} - \frac{\alpha^2}{(\alpha+\beta)^2}\\
&= \frac{\alpha\big[(\alpha+1)(\alpha+\beta) - \alpha(\alpha+\beta+1)\big]}{(\alpha+\beta)^2(\alpha+\beta+1)}
= \frac{\alpha\,\new{\beta}}{(\alpha+\beta)^2(\alpha+\beta+1)}.
\end{align*}
\end{enumerate}
These are the prior moments quoted above. Replacing $\alpha \to y + \alpha$ and $\beta \to \beta + n - y$ gives the posterior moments.
\end{homework}

\paragraph{Posterior predictive.}
The posterior predictive distribution is
\[
p(\tilde y \mid y) = \int p(\tilde y \mid \theta)\,p(\theta \mid y)\dd\theta.
\]
For $\tilde n \ge 1$ new trials,
\begin{align*}
p(\tilde y \mid y) &= \int \operatorname{Bin}(\tilde y \mid \tilde n, \theta)\,\operatorname{Beta}(\theta \mid y + \alpha, \beta + n - y)\dd\theta\\
&= \new{\binom{\tilde n}{\tilde y}}\new{\frac{1}{B(y+\alpha, \beta+n-y)}}\int \new{\theta^{\tilde y}(1-\theta)^{\tilde n - \tilde y}}\,\new{\theta^{y+\alpha-1}(1-\theta)^{\beta+n-y-1}}\dd\theta\\
&= \binom{\tilde n}{\tilde y}\frac{\new{B(\tilde y + y + \alpha,\ \beta + n + \tilde n - y - \tilde y)}}{B(y + \alpha,\ \beta + n - y)}.
\end{align*}

\begin{example}
Let us consider that we observed $y = 4$ out of $n = 5$. Also, let us choose $\alpha = 1$ and $\beta = 1$, so that $\operatorname{Beta}(\theta \mid 1, 1) = \operatorname{Uniform}(0,1)$. In that case,
\[
p(\tilde y \mid y = 4) = \binom{\tilde n}{\tilde y}\frac{B(\tilde y + 5,\ \tilde n - \tilde y + 2)}{B(5, 2)}.
\]
\end{example}


\begin{homework}{4}
Let $y_1, \dots, y_n$ be i.i.d.\ Poisson with parameter $\theta > 0$,
\[
p(y_i \mid \theta) = \frac{\theta^{y_i}e^{-\theta}}{y_i!}, \qquad y_i = 0, 1, 2, \dots,
\]
and let the prior be a Gamma density with shape $\alpha > 0$ and rate $\beta > 0$,
\[
p(\theta) = \frac{\beta^{\alpha}}{\Gamma(\alpha)}\theta^{\alpha-1}e^{-\beta\theta}, \qquad \theta > 0.
\]
\begin{enumerate}[label=(\roman*)]
\item Show that $p(\theta)$ is a conjugate prior for the likelihood $p(y_1, \dots, y_n \mid \theta)$.
\item Find the mean and variance of the posterior $p(\theta \mid y_1, \dots, y_n)$.
\item Find the posterior predictive distribution $p(\tilde y \mid y_1, \dots, y_n)$ of a new observation $\tilde y$.
\item Find the mean and variance of $p(\tilde y \mid y_1, \dots, y_n)$.
\end{enumerate}

\soln
\begin{enumerate}[label=(\roman*)]
\item Write $s = \sum_{i=1}^{n} y_i$. By independence,
\[
p(y_1, \dots, y_n \mid \theta) = \prod_{i=1}^{n}\frac{\theta^{y_i}e^{-\theta}}{y_i!} \propto \theta^{\new{s}}e^{\new{-n\theta}},
\]
and therefore
\[
p(\theta \mid y) \propto \theta^{s}e^{-n\theta}\cdot\theta^{\alpha-1}e^{-\beta\theta} = \theta^{\new{(\alpha+s)}-1}e^{-\new{(\beta+n)}\theta}.
\]
This is the kernel of a Gamma density, so $p(\theta \mid y) = \operatorname{Gamma}(\theta \mid \alpha_n, \beta_n)$ with
\[
\alpha_n = \alpha + \sum_{i=1}^{n} y_i, \qquad \beta_n = \beta + n.
\]
The posterior is in the same family as the prior, so the Gamma prior is conjugate to the Poisson likelihood.
\item For a $\operatorname{Gamma}(a, b)$ density the mean is $a/b$ and the variance $a/b^2$, so
\[
\E[\theta \mid y] = \frac{\alpha + \sum_i y_i}{\beta + n}, \qquad
\Var(\theta \mid y) = \frac{\alpha + \sum_i y_i}{(\beta + n)^2}.
\]
\item Integrating the Poisson likelihood of $\tilde y$ against the posterior,
\begin{align*}
p(\tilde y \mid y) &= \int_0^\infty \frac{\theta^{\tilde y}e^{-\theta}}{\tilde y!}\,\frac{\beta_n^{\alpha_n}}{\Gamma(\alpha_n)}\theta^{\alpha_n-1}e^{-\beta_n\theta}\dd\theta
= \frac{\beta_n^{\alpha_n}}{\Gamma(\alpha_n)\,\tilde y!}\int_0^\infty \theta^{\new{\alpha_n+\tilde y}-1}e^{-\new{(\beta_n+1)}\theta}\dd\theta\\
&= \frac{\beta_n^{\alpha_n}}{\Gamma(\alpha_n)\,\tilde y!}\,\new{\frac{\Gamma(\alpha_n+\tilde y)}{(\beta_n+1)^{\alpha_n+\tilde y}}}
= \frac{\Gamma(\alpha_n+\tilde y)}{\Gamma(\alpha_n)\,\tilde y!}\left(\frac{\beta_n}{\beta_n+1}\right)^{\alpha_n}\left(\frac{1}{\beta_n+1}\right)^{\tilde y},
\end{align*}
for $\tilde y = 0, 1, 2, \dots$, using $\int_0^\infty \theta^{a-1}e^{-b\theta}\dd\theta = \Gamma(a)/b^{a}$. This is a negative binomial distribution, $\tilde y \mid y \sim \operatorname{NegBin}\big(\alpha_n,\ \beta_n/(\beta_n+1)\big)$.
\item By the laws of total expectation and total variance, with $\E[\tilde y \mid \theta] = \Var(\tilde y \mid \theta) = \theta$,
\begin{align*}
\E[\tilde y \mid y] &= \E\big[\E[\tilde y \mid \theta] \,\big|\, y\big] = \E[\theta \mid y] = \frac{\alpha_n}{\beta_n},\\
\Var(\tilde y \mid y) &= \E\big[\Var(\tilde y \mid \theta) \,\big|\, y\big] + \Var\big(\E[\tilde y \mid \theta] \,\big|\, y\big)
= \frac{\alpha_n}{\beta_n} + \frac{\alpha_n}{\beta_n^2} = \new{\frac{\alpha_n(\beta_n+1)}{\beta_n^2}}.
\end{align*}
The predictive variance is larger than its mean, unlike a Poisson: the extra $\alpha_n/\beta_n^2$ is our remaining uncertainty about $\theta$.
\end{enumerate}
\end{homework}

%-----------------------------------------------------
\subsection{The Gaussian--Gaussian model: known variance}
Let us consider the likelihood for a single observation $y_i$:
\[
p(y_i \mid \theta) = \frac{1}{\sqrt{2\pi\sigma^2}}\exp\left(-\frac{1}{2\sigma^2}(y_i - \theta)^2\right),
\]
and assuming that $y_1, \dots, y_n$ are conditionally independent given $\theta$, then
\begin{align*}
p(y \mid \theta) &= \prod_{i=1}^{n} p(y_i \mid \theta) = \new{\frac{1}{(2\pi\sigma^2)^{n/2}}}\exp\left(-\frac{1}{2\sigma^2}\new{\sum_{i=1}^{n}}(y_i - \theta)^2\right)
\propto \exp\left(-\frac{1}{2\sigma^2}\sum_{i=1}^{n}(y_i - \theta)^2\right)\\
&= \exp\left(-\frac{1}{2\sigma^2}\sum_{i=1}^{n}\new{(y_i^2 - 2\theta y_i + \theta^2)}\right)
= \exp\left(-\frac{1}{2\sigma^2}\sum_{i=1}^{n}y_i^2 + \frac{2\theta}{2\sigma^2}\,\new{n\,\frac1n}\sum_{i=1}^{n}y_i - \frac{1}{2\sigma^2}\theta^2 \new{n}\right)\\
&= \exp\left(-\frac{1}{2\sigma^2}\sum_{i} y_i^2\right)\times\exp\left(\frac{\theta n\new{\bar y}}{\sigma^2} - \frac{1}{2\sigma^2}\theta^2 n\right)
\propto \exp\left(\new{\frac{n}{2\sigma^2}}\left(2\theta\bar y - \theta^2\right)\right)\\
&= \exp\left(-\frac{n}{2\sigma^2}\left(\theta^2 - 2\theta\bar y \new{{}+ \bar y^2 - \bar y^2}\right)\right)
= \exp\left(-\frac{n}{2\sigma^2}\new{(\theta - \bar y)^2} + \frac{\bar y^2 n}{2\sigma^2}\right)\\
&= \exp\left(-\frac{n}{2\sigma^2}(\theta - \bar y)^2\right)\new{\exp\left(\frac{n\bar y^2}{2\sigma^2}\right)}
\propto \exp\left(-\frac{n}{2\sigma^2}(\theta - \bar y)^2\right).
\end{align*}
So
\[
p(y \mid \theta) \propto \exp\left(-\frac{n}{2\sigma^2}(\theta - \bar y)^2\right).
\]
Also, consider a Gaussian prior
\[
p(\theta) = \frac{1}{(2\pi\tau_0^2)^{1/2}}\exp\left(-\frac{1}{2\tau_0^2}(\theta - \mu_0)^2\right),
\qquad\text{which}\qquad
p(\theta) \propto \exp\left(-\frac{1}{2\tau_0^2}(\theta - \mu_0)^2\right).
\]
Then
\[
p(y \mid \theta)\,p(\theta) \propto \exp\left[-\frac{n}{2\sigma^2}(\theta - \bar y)^2 - \frac{1}{2\tau_0^2}(\theta - \mu_0)^2\right].
\]
The exponent is
\begin{align*}
&-\frac{n}{2\sigma^2}(\theta^2 - 2\theta\bar y + \bar y^2) - \frac{1}{2\tau_0^2}(\theta^2 - 2\theta\mu_0 + \mu_0^2)\\
&= -\frac12\left[\theta^2\new{\left(\frac{n}{\sigma^2} + \frac{1}{\tau_0^2}\right)} - 2\theta\new{\left(\frac{\bar y n}{\sigma^2} + \frac{\mu_0}{\tau_0^2}\right)} + \bar y^2\frac{n}{\sigma^2} + \frac{\mu_0^2}{\tau_0^2}\right].
\end{align*}
The bracket equals
\[
\underbrace{\left(\frac{n}{\sigma^2} + \frac{1}{\tau_0^2}\right)}_{\new{1/\tau_n^2}}
\Bigg[\theta^2 - 2\theta\underbrace{\frac{\frac{\mu_0}{\tau_0^2} + \frac{\bar y n}{\sigma^2}}{\new{\frac{n}{\sigma^2} + \frac{1}{\tau_0^2}}}}_{\new{\mu_n}}\Bigg]
+ \underbrace{\left(\bar y^2\frac{n}{\sigma^2} + \frac{\mu_0^2}{\tau_0^2}\right)}_{\new{C}}.
\]
Bringing back the $-\frac12$:
\[
-\frac{1}{2\tau_n^2}\left(\theta^2 - 2\theta\mu_n \new{{}+ \mu_n^2}\right) \new{{}+ \frac{\mu_n^2}{2\tau_n^2}} - \frac12 C
= -\frac{1}{2\tau_n^2}\new{(\theta - \mu_n)^2} + \new{K}.
\]
Then
\[
p(y \mid \theta)\,p(\theta) \propto \exp\left(-\frac{1}{2\tau_n^2}(\theta - \mu_n)^2\right)\times\exp(K)
\ \Rightarrow\
p(y \mid \theta)\,p(\theta) \propto \exp\left[-\frac{1}{2\tau_n^2}(\theta - \mu_n)^2\right].
\]
Since $p(\theta \mid y) \propto p(y \mid \theta)\,p(\theta)$,
\[
p(\theta \mid y) \propto \exp\left[-\frac{1}{2\tau_n^2}(\theta - \mu_n)^2\right],
\qquad
\theta \mid y \sim \Normal(\mu_n, \tau_n^2),
\]
with
\[
\frac{1}{\tau_n^2} = \frac{n}{\sigma^2} + \frac{1}{\tau_0^2},
\qquad
\mu_n = \frac{\frac{\mu_0}{\tau_0^2} + \frac{\bar y n}{\sigma^2}}{\frac{n}{\sigma^2} + \frac{1}{\tau_0^2}}.
\]
Thus
\[
\E[\theta \mid y] = \mu_n, \qquad \Var(\theta \mid y) = \tau_n^2.
\]

\paragraph{Posterior predictive distribution.}
\[
p(\tilde y \mid y) = \int p(\tilde y \mid \theta, y)\,p(\theta \mid y)\dd\theta = \int p(\tilde y \mid \theta)\,p(\theta \mid y)\dd\theta.
\]
With $m$ the new number of observations and $\bar{\tilde y}$ their mean, the same split $\sum_i(\tilde y_i - \theta)^2 = \sum_i(\tilde y_i - \bar{\tilde y})^2 + m(\bar{\tilde y} - \theta)^2$ used for the likelihood gives
\[
p(\tilde y \mid y) \propto \exp\left(-\frac{1}{2\sigma^2}\sum_{i=1}^{m}(\tilde y_i - \bar{\tilde y})^2\right)\int \mathrm{d}\theta\,\exp\left(-\frac{m}{2\sigma^2}(\theta - \bar{\tilde y})^2\right)\exp\left(-\frac{1}{2\tau_n^2}(\theta - \mu_n)^2\right).
\]
The first factor depends on $\tilde y$, so it must be kept: we want a distribution over $\tilde y$, and only factors free of $\tilde y$ can go into the proportionality constant. Every factor is the exponential of a quadratic in $(\tilde y, \theta)$, and integrating $\theta$ out of a joint Gaussian leaves a Gaussian, so $p(\tilde y \mid y)$ is a (multivariate) Normal. Its moments are, with $\mathbf{1}$ the vector of $m$ ones,
\[
\tilde y \mid y \sim \Normal\left(\mu_n\mathbf{1},\ \sigma^2 I_m + \tau_n^2\,\mathbf{1}\mathbf{1}^{\top}\right),
\]
since $\tilde y_i = \theta + \varepsilon_i$ with $\theta \mid y \sim \Normal(\mu_n, \tau_n^2)$ shared by all new observations and $\varepsilon_i \sim \Normal(0, \sigma^2)$ independent. Also,
\[
\E[\tilde y \mid y] = \E\big[\E[\tilde y \mid \theta, y] \mid y\big],
\]
\[
\E[\tilde y \mid y] = \int \mathrm{d}\tilde y\,\tilde y\int p(\tilde y \mid \theta)\,p(\theta \mid y)\dd\theta
= \int \mathrm{d}\theta\,p(\theta \mid y)\underbrace{\int \mathrm{d}\tilde y\,\tilde y\,p(\tilde y \mid \theta)}_{\theta}
= \mu_n.
\]
For $m = 1$,
\[
\Var(\tilde y \mid y) = \sigma^2 + \tau_n^2.
\]

%-----------------------------------------------------
\subsection{Known mean (unknown variance)}
We saw before that for iid $y_1, \dots, y_n$
\[
p(y \mid \theta) = \prod_{i=1}^{n} p(y_i \mid \theta) = \frac{1}{(2\pi\sigma^2)^{n/2}}\exp\left(-\frac{1}{2\sigma^2}\sum_{i=1}^{n}(y_i - \theta)^2\right),
\]
where $p(y_i \mid \theta) = \frac{1}{\sqrt{2\pi}\,\sigma}\exp\left(-\frac{1}{2\sigma^2}(y_i - \theta)^2\right)$.

But now we consider that $\theta$ is known and $\sigma^2$ is unknown. Thus, we can rewrite $p(y \mid \theta)$ as
\[
p(y \mid \sigma^2) = \frac{1}{(2\pi\sigma^2)^{n/2}}\exp\left(-\frac{1}{2\sigma^2}\sum_{i=1}^{n}(y_i - \theta)^2\right)
\propto \frac{1}{(\sigma^2)^{n/2}}\exp\left(-\frac{nv}{2\sigma^2}\right),
\qquad
v = \frac1n\sum_{i=1}^{n}(y_i - \theta)^2.
\]
The appropriate prior for this is the \emph{inverse-gamma distribution}
\[
p(\phi) = \frac{\beta^\alpha}{\Gamma(\alpha)}\phi^{-(\alpha+1)}e^{-\beta/\phi}, \qquad \phi > 0,
\]
with
\[
\E[\phi] = \frac{\beta}{\alpha - 1}\ (\alpha > 1), \qquad \Var(\phi) = \frac{\beta^2}{(\alpha-1)^2(\alpha-2)}\ (\alpha > 2).
\]
Setting $\phi = \sigma^2$ we obtain
\[
p(\sigma^2) \propto (\sigma^2)^{-(\alpha+1)}e^{-\beta/\sigma^2}.
\]
Thus, the posterior is
\begin{align*}
p(\sigma^2 \mid y) &\propto p(y \mid \sigma^2)\,p(\sigma^2)
\propto (\sigma^2)^{\new{-n/2 - (\alpha+1)}}\exp\left[-\left(\new{\frac{nv}{2\sigma^2} + \frac{\beta}{\sigma^2}}\right)\right]\\
&= (\sigma^2)^{-(\new{\frac n2 + \alpha} + 1)}\exp\left[-\new{\frac{1}{\sigma^2}}\left(\new{\frac{nv}{2} + \beta}\right)\right],
\end{align*}
\[
p(\sigma^2 \mid y) = \frac{\beta'^{\alpha'}}{\Gamma(\alpha')}(\sigma^2)^{-(\new{\alpha'}+1)}\exp\left[-\frac{1}{\sigma^2}\new{\beta'}\right],
\qquad
\alpha' = \frac n2 + \alpha, \quad \beta' = \frac{nv}{2} + \beta.
\]

\paragraph{Posterior predictive distribution.}
\[
p(\tilde y \mid y) = \int \mathrm{d}\sigma^2\,p(\tilde y \mid \sigma^2, \theta)\,p(\sigma^2 \mid y),
\qquad
p(\tilde y \mid \sigma^2, \theta) = \frac{1}{(2\pi\sigma^2)^{m/2}}\exp\left[-\frac{1}{2\sigma^2}\sum_{i=1}^{m}(\tilde y_i - \theta)^2\right].
\]
Thus
\begin{align*}
p(\tilde y \mid y) &\propto \int \mathrm{d}\sigma^2\,\frac{1}{(\sigma^2)^{m/2}}\exp\left[-\frac{1}{2\sigma^2}\sum_{i=1}^{m}(\tilde y_i - \theta)^2\right](\sigma^2)^{-(\alpha'+1)}\exp\left[-\frac{1}{\sigma^2}\beta'\right]\\
&= \int \mathrm{d}\sigma^2\,(\sigma^2)^{-(\new{\frac m2 + \alpha'} + 1)}\exp\left[-\frac{1}{\sigma^2}\left(\new{\beta' + \frac12\sum_{i=1}^{m}(\tilde y_i - \theta)^2}\right)\right]
= \int \mathrm{d}\sigma^2\,(\sigma^2)^{-(\new{a}+1)}\exp\left[-\frac{\new{b}}{\sigma^2}\right],
\end{align*}
with
\[
a = \frac m2 + \alpha', \qquad b = \beta' + \frac12\sum_{i=1}^{m}(\tilde y_i - \theta)^2.
\]
Recall that $p(\phi) = \frac{\beta^\alpha}{\Gamma(\alpha)}\phi^{-(\alpha+1)}e^{-\beta/\phi}$; therefore
\[
\int \mathrm{d}\phi\,p(\phi) = \frac{\beta^\alpha}{\Gamma(\alpha)}\int \phi^{-(\alpha+1)}e^{-\beta/\phi}\dd\phi = 1
\quad\text{and thus}\quad
\int \phi^{-(\alpha+1)}e^{-\beta/\phi}\dd\phi = \frac{\Gamma(\alpha)}{\beta^\alpha}.
\]
Thus, keeping the constants,
\[
p(\tilde y \mid y) = \frac{1}{(2\pi)^{m/2}}\frac{\beta'^{\alpha'}}{\Gamma(\alpha')}\int \mathrm{d}\sigma^2\,(\sigma^2)^{-(a+1)}e^{-b/\sigma^2}
= \frac{1}{(2\pi)^{m/2}}\frac{\beta'^{\alpha'}}{\Gamma(\alpha')}\frac{\Gamma(a)}{b^a},
\]
\[
p(\tilde y \mid y) = \frac{1}{(2\pi)^{m/2}}\frac{(\beta')^{\alpha'}\,\Gamma(\frac m2 + \alpha')}{\Gamma(\alpha')\left[\beta' + \frac12\sum_{i=1}^{m}(\tilde y_i - \theta)^2\right]^{\frac m2 + \alpha'}}.
\]
If $m = 1$:
\[
p(\tilde y \mid y) = \frac{1}{(2\pi)^{1/2}}\frac{(\beta')^{\alpha'}\,\Gamma(\frac12 + \alpha')}{\Gamma(\alpha')\left[\beta' + \frac12(\tilde y_1 - \theta)^2\right]^{\frac12 + \alpha'}}.
\]
This is a $t$-Student distribution with $2\alpha'$ degrees of freedom, with
\[
\E[\tilde y \mid y] = \theta, \qquad \Var[\tilde y \mid y] = \frac{\beta'}{\alpha' - 1}.
\]

%-----------------------------------------------------
\subsection{Unknown $\mu$ and $\sigma^2$}
Now, we want to make inference about $\mu$ and $\sigma^2$. This is, $\theta = (\mu, \sigma^2)$ and
\[
p(\mu, \sigma^2 \mid y) \propto p(y \mid \mu, \sigma^2)\,p(\mu, \sigma^2).
\]
The trick here is to write
\[
p(\mu, \sigma^2) = p(\mu \mid \sigma^2)\,p(\sigma^2)
\]
and choose
\[
\mu \mid \sigma^2 \sim \Normal(\mu_0, \sigma^2/\kappa_0), \qquad \sigma^2 \sim \text{Inv-}\chi^2(\nu_0, \sigma_0^2).
\]
This is
\[
p(\mu \mid \sigma^2) = \left(\frac{\kappa_0}{2\pi\sigma^2}\right)^{1/2}\exp\left(-\frac{\kappa_0(\mu - \mu_0)^2}{2\sigma^2}\right),
\qquad
p(\sigma^2) \propto (\sigma^2)^{-(\nu_0/2 + 1)}\exp\left(-\frac{\nu_0\sigma_0^2}{2\sigma^2}\right),
\]
knowing that
\[
\int_0^\infty x^{-(a+1)}e^{-b/x}\dd x = \frac{\Gamma(a)}{b^a}.
\]
Together,
\[
p(\mu, \sigma^2) = \text{N-Inv-}\chi^2\left(\mu, \sigma^2 \,\Big|\, \mu_0, \frac{\sigma_0^2}{\kappa_0}; \nu_0, \sigma_0^2\right).
\]

\paragraph{Likelihood.} With $\theta = (\mu, \sigma^2)$,
\[
p(y \mid \mu, \sigma^2) = \frac{1}{(2\pi\sigma^2)^{n/2}}\exp\left(-\frac{1}{2\sigma^2}\sum_{i=1}^{n}(y_i - \mu)^2\right)
= \frac{1}{(2\pi\sigma^2)^{n/2}}\exp\left(-\frac{1}{2\sigma^2}\new{\left[(n-1)s^2 + n(\bar y - \mu)^2\right]}\right),
\]
where
\[
s^2 = \frac{1}{n-1}\sum_{i=1}^{n}(y_i - \bar y)^2
\]
is the sampling variance.

\paragraph{Posterior.} $p(\mu, \sigma^2 \mid y) \propto p(y \mid \mu, \sigma^2)\,p(\mu, \sigma^2)$, giving
\[
p(\mu, \sigma^2 \mid y) = \text{N-Inv-}\chi^2\left(\mu, \sigma^2 \,\Big|\, \mu_n, \frac{\sigma_n^2}{\kappa_n}; \nu_n, \sigma_n^2\right),
\]
with
\begin{align*}
\mu_n &= \left(\frac{\kappa_0}{\kappa_0 + n}\right)\mu_0 + \frac{n}{\kappa_0 + n}\bar y \qquad\text{(weighted average)},\\
\kappa_n &= \kappa_0 + n,\\
\nu_n &= \nu_0 + n,\\
\nu_n\sigma_n^2 &= \nu_0\sigma_0^2 + (n-1)s^2 + \left(\frac{\kappa_0 n}{\kappa_0 + n}\right)(\bar y - \mu_0)^2.
\end{align*}
Then,
\[
\mu \mid \sigma^2, y \sim \Normal(\mu_n, \sigma^2/\kappa_n), \qquad \sigma^2 \mid y \sim \text{Inv-}\chi^2(\nu_n, \sigma_n^2).
\]

\paragraph{Posterior predictive.}
\[
p(\tilde y) = \int \mathrm{d}\theta\,p(\tilde y \mid \theta)\,p(\theta),
\qquad
p(\tilde y \mid y) = \int \mathrm{d}\theta\,p(\tilde y \mid \theta, y)\,p(\theta \mid y) = \int \mathrm{d}\theta\,\new{p(\tilde y \mid \theta)}\,p(\theta \mid y),
\]
\[
p(\tilde y \mid y) = \new{\int \mathrm{d}\sigma^2\int \mathrm{d}\mu}\,p(\tilde y \mid \new{\mu, \sigma^2})\,p(\new{\mu, \sigma^2} \mid y).
\]
Since $\int \tilde y\,p(\tilde y \mid \mu, \sigma^2)\dd\tilde y = \mu$,
\[
\int \tilde y\,p(\tilde y \mid y)\dd\tilde y = \int \mu\,p(\mu \mid y)\dd\mu.
\]

\paragraph{Marginal posterior of $\mu$.}
From the likelihood and the prior,
\begin{align*}
p(y \mid \mu, \sigma^2) &= \frac{1}{(2\pi\sigma^2)^{n/2}}\exp\left(-\frac{1}{2\sigma^2}\left[(n-1)s^2 + n(\bar y - \mu)^2\right]\right),\\
p(\mu, \sigma^2) &\propto \left(\frac{\kappa_0}{2\pi\sigma^2}\right)^{1/2}\exp\left(-\frac{\kappa_0(\mu - \mu_0)^2}{2\sigma^2}\right)\times(\sigma^2)^{-(\nu_0/2 + 1)}\exp\left(-\frac{\nu_0\sigma_0^2}{2\sigma^2}\right),
\end{align*}
we get
\[
p(\mu, \sigma^2 \mid y) \propto \frac{1}{(\sigma^2)^{n/2}}\frac{1}{(\sigma^2)^{1/2}}\frac{1}{(\sigma^2)^{(\nu_0/2 + 1)}}\exp\left(-\frac{\new{B}}{\sigma^2}\right)
= (\sigma^2)^{\new{-(\frac{\nu_0}{2} + \frac n2 + \frac12 + 1)}}\exp\left(-\frac{B}{\sigma^2}\right),
\]
\[
p(\mu, \sigma^2 \mid y) = \new{K}(\sigma^2)^{-(\new{A}+1)}\exp\left(-\frac{B}{\sigma^2}\right),
\]
with
\[
A = \frac{\nu_0 + n + 1}{2} = \frac{\nu_n + 1}{2},
\qquad
B = \frac12\left[\nu_n\sigma_n^2 + \kappa_n(\mu - \mu_n)^2\right].
\]
Then
\[
\int p(\mu, \sigma^2 \mid y)\dd\sigma^2 = p(\mu \mid y),
\]
\[
\int p(\mu, \sigma^2 \mid y)\dd\sigma^2 = \new{K}\int \new{(\sigma^2)^{-(A+1)}e^{-B/\sigma^2}}\dd\sigma^2 = K\new{\frac{\Gamma(A)}{B^A}} = K\,\Gamma(A)\,\new{B^{-A}},
\]
so
\[
p(\mu \mid y) \propto \new{\Gamma\left(\frac{\nu_n + 1}{2}\right)}\new{\left[1 + \frac{\kappa_n(\mu - \mu_n)^2}{\nu_n\sigma_n^2}\right]^{-\frac{\nu_n + 1}{2}}}.
\]

%=====================================================
\section{Exponential family}
For $n$ observations, a distribution that can be written as
\[
p(\vec y \mid \vec\theta) \propto g(\vec\theta)^{n}\exp\left[f(\vec\theta)^{\top}T(\vec y)\right]
\]
is said to belong to the \emph{exponential family}. $T(\vec y)$ is the \emph{sufficient statistic} for $\vec\theta$, and $f(\vec\theta)$ is the \emph{natural parameter}.

If the prior can be written as
\[
p(\vec\theta) \propto g(\vec\theta)^{\eta}\exp\left[f(\vec\theta)^{\top}\vec\nu\right],
\]
then the posterior is
\[
p(\vec\theta \mid \vec y) \propto g(\vec\theta)^{n+\eta}\exp\left[f(\vec\theta)^{\top}\big(T(\vec y) + \vec\nu\big)\right]:
\]
same $f(\vec\theta)$ and same $g(\vec\theta)$.

\begin{example}[Beta--Binomial]
\[
p(y \mid \theta) = \binom{m}{y}\theta^{y}(1-\theta)^{m-y} = \binom{m}{y}\exp\left[\log\left(\theta^{y}(1-\theta)^{m-y}\right)\right]
= \binom{m}{y}\exp\left[y\log\theta + (m-y)\log(1-\theta)\right]
\]
\[
\propto \exp\left(f(\theta)^{\top}T(y)\right),
\qquad
f(\theta) = \begin{pmatrix}\log\theta\\ \log(1-\theta)\end{pmatrix},
\]
\[
T(y) = \begin{pmatrix}y\\ m-y\end{pmatrix}\ \text{(sufficient statistics vector)},
\quad
g(\theta) = 1.
\]
Beta prior:
\[
p(\theta) \propto \theta^{\alpha-1}(1-\theta)^{\beta-1} = \exp\left[(\alpha-1)\log\theta + (\beta-1)\log(1-\theta)\right].
\]
Thus
\[
\vec\nu = \begin{pmatrix}\alpha - 1\\ \beta - 1\end{pmatrix}.
\]
\end{example}


\begin{homework}{4}
\begin{enumerate}[label=(\roman*)]
\item For the Poisson likelihood $p(y_1, \dots, y_n \mid \theta)$ of the previous exercise, find the natural parameter $f(\theta)$, the function $g(\theta)$ and the sufficient statistic $T(y)$.
\item Using this notation, find a suitable prior for the Poisson likelihood and identify the resulting distribution.
\end{enumerate}
Make sure all distributions are normalized and the parameters are consistent with the exponential-family notation.

\soln
\begin{enumerate}[label=(\roman*)]
\item
\[
p(y \mid \theta) = \prod_{i=1}^{n}\frac{\theta^{y_i}e^{-\theta}}{y_i!}
= \Big(\prod_i \frac{1}{y_i!}\Big)\,e^{-n\theta}\,\theta^{\sum_i y_i}
\propto \new{\big(e^{-\theta}\big)^{n}}\exp\Big[\new{\log\theta}\,\new{\sum_i y_i}\Big],
\]
so
\[
g(\theta) = e^{-\theta}, \qquad f(\theta) = \log\theta, \qquad T(y) = \sum_{i=1}^{n} y_i.
\]
\item A conjugate prior has the form
\[
p(\theta) \propto g(\theta)^{\eta}\exp\big[f(\theta)\,\nu\big] = e^{-\eta\theta}\exp(\nu\log\theta) = \new{\theta^{\nu}e^{-\eta\theta}},
\]
the kernel of a $\operatorname{Gamma}(\nu + 1, \eta)$ density (shape $\nu + 1$, rate $\eta$). Normalized,
\[
p(\theta) = \frac{\eta^{\nu+1}}{\Gamma(\nu+1)}\theta^{\nu}e^{-\eta\theta}, \qquad \nu > -1,\ \eta > 0.
\]
The posterior is then
\[
p(\theta \mid y) \propto g(\theta)^{n+\eta}\exp\big[f(\theta)\big(T(y) + \nu\big)\big] = \theta^{\new{\nu + \sum_i y_i}}e^{-\new{(\eta+n)}\theta},
\]
that is $\operatorname{Gamma}\big(\nu + 1 + \sum_i y_i,\ \eta + n\big)$. With $\alpha = \nu + 1$ and $\beta = \eta$ this is the result of the previous exercise.
\end{enumerate}
\end{homework}

%=====================================================
\section{Multivariate normal model with known variance}
\paragraph{Likelihood.}
\[
p(\vec y \mid \vmu, \Sigma) = |\Sigma|^{-1/2}(2\pi)^{-D/2}\exp\left(-\frac12(\vec y - \vmu)^{\top}\Sigma^{-1}(\vec y - \vmu)\right),
\]
with $\Sigma$ the covariance matrix with entries
\[
\Sigma = \E\left[(\vec Y - \vmu)(\vec Y - \vmu)^{\top}\right],
\qquad
\Sigma_{jk} = \Cov(Y_j, Y_k) = \E(Y_jY_k) - \E(Y_j)\E(Y_k).
\]
For a population of size $N$, we can find $\Sigma$ as
\[
\Sigma = \frac1N\sum_{i=1}^{N}(\vec y_i - \vmu)(\vec y_i - \vmu)^{\top},
\qquad
\hat\mu = \frac1N\sum_{i=1}^{N}\vec y_i.
\]
For a sample of size $n$,
\[
\hat\Sigma = \frac{1}{n-1}\sum_{i=1}^{n}(\vec y_i - \hat\mu)(\vec y_i - \hat\mu)^{\top}.
\]
For $n$ iid observations $\vec y_1, \dots, \vec y_n \in \R^D$, the likelihood is
\[
p(\vec y_1, \dots, \vec y_n \mid \vmu, \Sigma) \propto |\Sigma|^{-n/2}\exp\left[-\frac12\sum_{i=1}^{n}(\vec y_i - \vmu)^{\top}\Sigma^{-1}(\vec y_i - \vmu)\right],
\qquad
\hat\mu = \frac1n\sum_{i=1}^{n}\vec y_i.
\]

To find $\Sigma^{-1}$ we can find the spectral decomposition of $\Sigma$ as
\[
\Sigma = \sum_{i=1}^{D}\lambda_i\vec u_i\vec u_i^{\top},
\qquad
\Sigma^{-1} = \sum_{j=1}^{D}\frac{1}{\lambda_j}\vec u_j\vec u_j^{\top}.
\]
Indeed,
\[
\Sigma\Sigma^{-1} = \sum_i\sum_j\frac{\lambda_i}{\lambda_j}\vec u_i\underbrace{\vec u_i^{\top}\vec u_j}_{\new{\delta_{ij}}}\vec u_j^{\top}
= \sum_j\frac{\new{\lambda_j}}{\lambda_j}\new{\vec u_j}\vec u_j^{\top} = \new{\one}.
\]

Notice that, in
\[
\sum_{i=1}^{n}(\vec y_i - \vmu)^{\top}\Sigma^{-1}(\vec y_i - \vmu),
\]
taking $\vec\zeta = \vec y_\ell - \vmu$ and $\vec\zeta = \sum_j\zeta_j\vec e_j$ with $\vec e_j$ the canonical basis, and writing $\Sigma^{-1} = \sum_{k\ell}\sigma_{k\ell}\vec e_k\vec e_\ell^{\top}$, then
\[
\Big(\sum_j\zeta_j\vec e_j^{\top}\Big)\Big(\sum_{k\ell}\sigma_{k\ell}\vec e_k\vec e_\ell^{\top}\Big)\Big(\sum_m\zeta_m\vec e_m\Big)
= \sum_{jk\ell m}\zeta_j\zeta_m\new{\delta_{jk}\delta_{\ell m}}\sigma_{k\ell}
= \sum_{k\ell}\new{\zeta_k\zeta_\ell}\sigma_{k\ell}
= \sum_k\zeta_k\new{\sum_\ell}\zeta_\ell\sigma_{k\ell}.
\]

\subsection{Prior for $\vmu$ with known $\Sigma$}
\[
\vmu \sim \Normal(\vmu_0, \Lambda_0),
\qquad
p(\vmu \mid \Lambda_0) \propto \exp\left[-\frac12(\vmu - \vmu_0)^{\top}\Lambda_0^{-1}(\vmu - \vmu_0)\right].
\]

\paragraph{Posterior.}
\begin{align*}
p(\vmu \mid \vec y_1, \dots, \vec y_n, \Sigma) &\propto \exp\left[-\frac12\sum_{i=1}^{n}(\vec y_i - \vmu)^{\top}\Sigma^{-1}(\vec y_i - \vmu)\right]\times\exp\left[-\frac12(\vmu - \vmu_0)^{\top}\Lambda_0^{-1}(\vmu - \vmu_0)\right]\\
&= \exp\left[-\frac12\left(\sum_{i=1}^{n}(\vec y_i - \vmu)^{\top}\Sigma^{-1}(\vec y_i - \vmu) \new{{}+ (\vmu - \vmu_0)^{\top}\Lambda_0^{-1}(\vmu - \vmu_0)}\right)\right].
\end{align*}
Expanding and dropping the terms without $\vmu$,
\[
\exp\Big[-\frac12\Big(n\vmu^{\top}\Sigma^{-1}\vmu - \sum_i\vec y_i^{\top}\Sigma^{-1}\vmu - \vmu^{\top}\Sigma^{-1}\sum_i\vec y_i + \vmu^{\top}\Lambda_0^{-1}\vmu - \vmu^{\top}\Lambda_0^{-1}\vmu_0 - \vmu_0^{\top}\Lambda_0^{-1}\vmu\Big)\Big]
\]
\[
\propto \exp\Big[-\frac12\Big(\vmu^{\top}\new{(n\Sigma^{-1} + \Lambda_0^{-1})}\vmu - \vmu^{\top}\new{\Big(\Sigma^{-1}\sum_{i=1}^{n}\vec y_i + \Lambda_0^{-1}\vmu_0\Big)} - \new{\Big(\sum_{i=1}^{n}\vec y_i^{\top}\Sigma^{-1} + \vmu_0^{\top}\Lambda_0^{-1}\Big)}\vmu\Big)\Big].
\]
The argument is
\[
\vmu^{\top}\new{\Lambda_n^{-1}}\vmu - \vmu^{\top}\left(\new{n\Sigma^{-1}\bar{\vec y}} + \Lambda_0^{-1}\vmu_0\right) - \left(\new{n\bar{\vec y}^{\top}}\Sigma^{-1} + \vmu_0^{\top}\Lambda_0^{-1}\right)\vmu,
\qquad\text{with}\quad \Lambda_n^{-1} = n\Sigma^{-1} + \Lambda_0^{-1}.
\]
Let us define
\[
\left(n\Sigma^{-1}\bar{\vec y} + \Lambda_0^{-1}\vmu_0\right) = \new{\Lambda_n^{-1}\vmu_n}
\ \Rightarrow\
\vmu_n = \new{\Lambda_n}\left(n\Sigma^{-1}\bar{\vec y} + \Lambda_0^{-1}\vmu_0\right),
\qquad
\left(n\bar{\vec y}^{\top}\Sigma^{-1} + \vmu_0^{\top}\Lambda_0^{-1}\right) = \vmu_n^{\top}\Lambda_n^{-1}.
\]
Then the argument is
\begin{align*}
&\vmu^{\top}\Lambda_n^{-1}\vmu - \vmu^{\top}\Lambda_n^{-1}\vmu_n - \vmu_n^{\top}\Lambda_n^{-1}\vmu \new{{}+ \vmu_n^{\top}\Lambda_n^{-1}\vmu_n - \vmu_n^{\top}\Lambda_n^{-1}\vmu_n}\\
&= \new{(\vmu^{\top} - \vmu_n^{\top})}\Lambda_n^{-1}\vmu - \new{(\vmu^{\top} - \vmu_n^{\top})}\Lambda_n^{-1}\vmu_n - \vmu_n^{\top}\Lambda_n^{-1}\vmu_n\\
&= (\vmu - \vmu_n)^{\top}\new{\left(\Lambda_n^{-1}\vmu - \Lambda_n^{-1}\vmu_n\right)} - \vmu_n^{\top}\Lambda_n^{-1}\vmu_n
= (\vmu - \vmu_n)^{\top}\new{\Lambda_n^{-1}(\vmu - \vmu_n)} - \vmu_n^{\top}\Lambda_n^{-1}\vmu_n.
\end{align*}
Thus
\[
p(\vmu \mid \vec y_1, \dots, \vec y_n, \Sigma) \propto \exp\left[-\frac12(\vmu - \vmu_n)^{\top}\Lambda_n^{-1}(\vmu - \vmu_n)\right],
\]
\[
p(\vmu \mid \vec y_1, \dots, \vec y_n, \Sigma) = (2\pi)^{-D/2}|\Lambda_n|^{-1/2}\exp\left[-\frac12(\vmu - \vmu_n)^{\top}\Lambda_n^{-1}(\vmu - \vmu_n)\right],
\]
with
\[
\Lambda_n^{-1} = n\Sigma^{-1} + \Lambda_0^{-1}
\qquad\text{and}\qquad
\vmu_n = \Lambda_n\left(n\Sigma^{-1}\bar{\vec y} + \Lambda_0^{-1}\vmu_0\right).
\]

\subsection{Posterior predictive distribution}
\[
p(\tilde y \mid y) = \int p(\tilde y \mid \theta)\,p(\theta \mid y)\dd\theta,
\qquad
p(\tilde{\vec y} \mid \vec y_1, \dots, \vec y_n, \Sigma) = \int p(\tilde{\vec y} \mid \vmu, \Sigma)\,p(\vmu \mid \vec y_1, \dots, \vec y_n, \Sigma)\dd\vmu,
\]
with
\[
p(\tilde{\vec y} \mid \vmu, \Sigma) \propto \exp\left(-\frac12(\tilde{\vec y} - \vmu)^{\top}\Sigma^{-1}(\tilde{\vec y} - \vmu)\right),
\]
\[
p(\vmu \mid \vec y_1, \dots, \vec y_n, \Sigma) \propto \exp\left[-\frac12(\vmu - \vmu_n)^{\top}\Lambda_n^{-1}(\vmu - \vmu_n)\right].
\]
\textbf{Mean:}
\[
\E[\tilde{\vec y} \mid \vec y_{1:n}] = \int \tilde{\vec y}\,p(\tilde{\vec y} \mid \vec y_{1:n}, \Sigma)\dd\tilde{\vec y}
= \int \mathrm{d}\vmu\,p(\vmu \mid \vec y_{1:n}, \Sigma)\underbrace{\int \mathrm{d}\tilde{\vec y}\,\tilde{\vec y}\,p(\tilde{\vec y} \mid \vmu, \Sigma)}_{\new{\vmu}}
= \int \mathrm{d}\vmu\,\new{\vmu}\,p(\vmu \mid \vec y_{1:n}, \Sigma) = \new{\vmu_n}.
\]
\textbf{Variance:}
\[
\Var(\tilde{\vec y} \mid \vec y_{1:n}) = \int\big(\tilde{\vec y} - \underbrace{\E(\tilde{\vec y})}_{\vmu_n}\big)\big(\tilde{\vec y} - \underbrace{\E(\tilde{\vec y})}_{\vmu_n}\big)^{\top}p(\tilde{\vec y} \mid \vec y_{1:n})\dd\tilde{\vec y}.
\]
Expanding $\tilde{\vec y}\tilde{\vec y}^{\top} - \tilde{\vec y}\vmu_n^{\top} - \vmu_n\tilde{\vec y}^{\top} + \vmu_n\vmu_n^{\top}$,
\begin{align*}
\Var(\tilde{\vec y} \mid \vec y_{1:n}) &= \int \tilde{\vec y}\tilde{\vec y}^{\top}p(\tilde{\vec y} \mid \vec y_{1:n})\dd\tilde{\vec y} - \underbrace{\Big(\int \tilde{\vec y}\,p(\tilde{\vec y} \mid \vec y_{1:n})\dd\tilde{\vec y}\Big)}_{\new{\vmu_n}}\vmu_n^{\top} - \vmu_n\underbrace{\Big(\int \tilde{\vec y}^{\top}p(\tilde{\vec y} \mid \vec y_{1:n})\dd\tilde{\vec y}\Big)}_{\new{\vmu_n^{\top}}} + \vmu_n\vmu_n^{\top}\\
&= \int \tilde{\vec y}\tilde{\vec y}^{\top}p(\tilde{\vec y} \mid \vec y_{1:n})\dd\tilde{\vec y} \new{{}- \vmu_n\vmu_n^{\top}}.
\end{align*}
Using $p(\tilde y \mid y) = \int \mathrm{d}\vmu\,p(\tilde{\vec y} \mid \vmu, \Sigma)\,p(\vmu \mid \vec y_{1:n})$,
\[
\int \tilde{\vec y}\tilde{\vec y}^{\top}p(\tilde{\vec y} \mid \vec y_{1:n})\dd\tilde{\vec y}
= \int \mathrm{d}\vmu\,p(\vmu \mid \vec y_{1:n})\underbrace{\Big(\int \mathrm{d}\tilde{\vec y}\,\tilde{\vec y}\tilde{\vec y}^{\top}p(\tilde{\vec y} \mid \vmu, \Sigma)\Big)}_{\new{\Sigma + \vmu\vmu^{\top}}}
= \new{\Sigma} + \underbrace{\int \new{\vmu\vmu^{\top}}p(\vmu \mid \vec y_{1:n})\dd\vmu}_{\new{\Lambda_n + \vmu_n\vmu_n^{\top}}}.
\]
Therefore
\[
\boxed{\Var(\tilde{\vec y} \mid \vec y_{1:n}) = \Sigma + \Lambda_n.}
\]

%=====================================================
\section{Linear Gaussian systems}
Let us consider that we have noisy measurements of an unknown vector $\vec z$ such that
\[
\vec y = W\vec z + \vec b + \vec\varepsilon,
\]
where $\vec\varepsilon \sim \Normal(\vec 0, \Sigma_y)$, such that
\[
\E[\vec y \mid \vec z] = W\vec z + \vec b + \underbrace{\E(\vec\varepsilon)}_{0}.
\]
Thus,
\[
p(\vec y \mid \vec z) = \Normal(\vec y \mid W\vec z + \vec b, \Sigma_y).
\]
For this model, we choose a prior
\[
p(\vec z) = \Normal(\vec z \mid \vmu_z, \Sigma_z).
\]
Thus, the posterior distribution is
\[
p(\vec z \mid \vec y) \propto p(\vec y \mid \vec z)\,p(\vec z).
\]
For $n$ observations, $\vec y_i = W_i\vec z + \vec b_i + \vec\varepsilon_i$, we have
\[
p(\vec z \mid \vec y_{1:n}) \propto p(\vec y_{1:n} \mid \vec z)\,p(\vec z),
\qquad
p(\vec y_{1:n} \mid \vec z) = \prod_{i=1}^{n}p(\vec y_i \mid \vec z),
\]
where
\[
p(\vec y_i \mid \vec z) \propto \exp\left[-\frac12(\vec y_i - W_i\vec z - \vec b_i)^{\top}\Sigma_y^{-1}(\vec y_i - W_i\vec z - \vec b_i)\right],
\qquad
p(\vec z) \propto \exp\left[-\frac12(\vec z - \vmu_z)^{\top}\Sigma_z^{-1}(\vec z - \vmu_z)\right].
\]
Then, in the exponent of $p(\vec z \mid \vec y_{1:n})$ (times $-2$) we have
\begin{align*}
\sum_{i=1}^{n}\Big(&\vec y_i^{\top}\Sigma_y^{-1}\vec y_i - \vec y_i^{\top}\Sigma_y^{-1}W_i\vec z - \vec y_i^{\top}\Sigma_y^{-1}\vec b_i
- \vec z^{\top}W_i^{\top}\Sigma_y^{-1}\vec y_i + \vec z^{\top}W_i^{\top}\Sigma_y^{-1}W_i\vec z + \vec z^{\top}W_i^{\top}\Sigma_y^{-1}\vec b_i\\
&- \vec b_i^{\top}\Sigma_y^{-1}\vec y_i + \vec b_i^{\top}\Sigma_y^{-1}W_i\vec z + \vec b_i^{\top}\Sigma_y^{-1}\vec b_i\Big)
+ \vec z^{\top}\Sigma_z^{-1}\vec z - \vec z^{\top}\Sigma_z^{-1}\vmu_z - \vmu_z^{\top}\Sigma_z^{-1}\vec z + \vmu_z^{\top}\Sigma_z^{-1}\vmu_z.
\end{align*}
I only need terms with $\vec z$. Thus
\[
\vec z^{\top}\new{\Big(\Sigma_z^{-1} + \sum_i W_i^{\top}\Sigma_y^{-1}W_i\Big)}\vec z
+ \new{2}\vec z^{\top}\sum_{i=1}^{n}\big(W_i^{\top}\Sigma_y^{-1}\vec b_i\big) - \new{2}\vec z^{\top}\sum_{i=1}^{n}\big(W_i^{\top}\Sigma_y^{-1}\vec y_i\big) - \new{2}\vec z^{\top}\Sigma_z^{-1}\vmu_z.
\]
Defining
\[
\Sigma_{z|y}^{-1} = \Sigma_z^{-1} + \sum_{i=1}^{n}\big(W_i^{\top}\Sigma_y^{-1}W_i\big),
\]
then
\[
\vec z^{\top}\new{\Sigma_{z|y}^{-1}}\vec z - 2\vec z^{\top}\underbrace{\Big[\sum_{i=1}^{n}W_i^{\top}\Sigma_y^{-1}\new{(\vec y_i - \vec b_i)} + \Sigma_z^{-1}\vmu_z\Big]}_{\new{\vec r}}
= \vec z^{\top}\Sigma_{z|y}^{-1}\vec z - 2\vec z^{\top}\new{\vec r}
= \vec z^{\top}\Sigma_{z|y}^{-1}\vec z - 2\vec z^{\top}\new{\Sigma_{z|y}^{-1}\vmu_{z|y}},
\]
with
\[
\vmu_{z|y} = \Sigma_{z|y}\Big(\sum_{i=1}^{n}W_i^{\top}\Sigma_y^{-1}(\vec y_i - \vec b_i) + \Sigma_z^{-1}\vmu_z\Big).
\]
Completing the square,
\begin{align*}
&\vec z^{\top}\Sigma_{z|y}^{-1}\vec z - \vec z^{\top}\Sigma_{z|y}^{-1}\vmu_{z|y} - \vmu_{z|y}^{\top}\Sigma_{z|y}^{-1}\vec z \new{{}+ \vmu_{z|y}^{\top}\Sigma_{z|y}^{-1}\vmu_{z|y} - \vmu_{z|y}^{\top}\Sigma_{z|y}^{-1}\vmu_{z|y}}\\
&= \new{(\vec z - \vmu_{z|y})^{\top}}\Sigma_{z|y}^{-1}\vec z - \new{(\vec z - \vmu_{z|y})^{\top}}\Sigma_{z|y}^{-1}\vmu_{z|y} - \vmu_{z|y}^{\top}\Sigma_{z|y}^{-1}\vmu_{z|y}\\
&= (\vec z - \vmu_{z|y})^{\top}\Sigma_{z|y}^{-1}\new{(\vec z - \vmu_{z|y})} - \vmu_{z|y}^{\top}\Sigma_{z|y}^{-1}\vmu_{z|y},
\end{align*}
so $p(\vec z \mid \vec y_{1:n}) = \Normal(\vec z \mid \vmu_{z|y}, \Sigma_{z|y})$.

%=====================================================
\section{Mixture of conjugate priors}
If $p(\theta \mid m = k)$ is a conjugate prior to a likelihood $p(y \mid \theta)$, such that
\[
p(\theta \mid y, m = k) = \frac{p(y \mid \theta, m = k)\,p(\theta \mid m = k)}{p(y \mid m = k)}
\]
for all $k \in \set{1, 2, \dots, K}$, then
\[
p(\theta) = \sum_k p(\theta \mid m = k)\,\underbrace{p(m = k)}_{\pi_k}, \qquad \sum_k \pi_k = 1,
\]
is also a conjugate prior to the likelihood $p(y \mid \theta)$.

The likelihood $p(y \mid \theta, m = k)$ is the same for all the priors $p(\theta \mid m = k)$. Thus,
\[
p(y \mid \theta, m = k) = p(y \mid \theta).
\]
Therefore
\[
p(\theta \mid y, m = k) = \frac{p(y \mid \theta)\,p(\theta \mid m = k)}{p(y \mid m = k)}
\qquad\text{and}\qquad
p(\theta \mid y, m = k)\,p(y \mid m = k) = p(y \mid \theta)\,p(\theta \mid m = k).
\]
Multiplying by $p(m = k)$ and summing up all the cases,
\[
\new{\sum_k p(m = k)}\,p(\theta \mid y, m = k)\,p(y \mid m = k) = p(y \mid \theta)\new{\sum_k p(m = k)}\,p(\theta \mid m = k).
\]
Taking $p(\theta) = \sum_k p(m = k)\,p(\theta \mid m = k)$ and $p(y \mid m = k)\,p(m = k) = p(m = k \mid y)\,p(y)$, we get
\[
\sum_k p(\theta \mid y, m = k)\,\new{p(m = k \mid y)} = \frac{p(y \mid \theta)\,\new{p(\theta)}}{\new{p(y)}} = \new{p(\theta \mid y)}.
\]

%=====================================================
\section{Model selection}
Let us suppose that we have different models, each labeled by $m \in \set{1, \dots, K}$, with likelihood $p(\mathcal{D} \mid \theta_m, m)$ and prior $p(\theta_m \mid m)$. Conditioning on $m$ specifies which model is being considered. Each model may have its own likelihood, parameters and prior.
\[
p(m \mid \mathcal{D}) = \frac{p(\mathcal{D} \mid m)\,p(m)}{\sum_{m'}p(\mathcal{D} \mid m')\,p(m')},
\]
and choose the model $\hat m$ such that
\[
\hat m = \arg\max_m\,p(m \mid \mathcal{D}).
\]
\[
p(\mathcal{D} \mid m) = \int \mathrm{d}\theta_m\,p(\mathcal{D} \mid \theta_m, m)\,p(\theta_m \mid m)
\]
is called the \emph{marginal likelihood} or \emph{model evidence}.

If the prior $p(m)$ is uniform, that is $p(m) = 1/K$, then
\[
\hat m = \arg\max_m\,p(\mathcal{D} \mid m).
\]

%=====================================================
\section*{References}
\begin{itemize}
\item Gelman, A., Carlin, J.~B., Stern, H.~S., Dunson, D.~B., Vehtari, A., and Rubin, D.~B. (2013). \emph{Bayesian Data Analysis} (3rd ed.). CRC Press. Chapters 1--3: Bayesian inference, single-parameter models (binomial, normal, Poisson), conjugacy and the exponential family, the normal model with unknown mean and variance, and the multivariate normal. Chapter 7: model comparison.
\item Hoff, P.~D. (2009). \emph{A First Course in Bayesian Statistical Methods}. Springer. Chapters 3, 5 and 7: binomial and Poisson models, the normal model, the multivariate normal model.
\item Murphy, K.~P. (2012). \emph{Machine Learning: A Probabilistic Perspective}. MIT Press. Chapter 4: Gaussian models, including linear Gaussian systems; Chapter 5: Bayesian statistics, mixtures of conjugate priors and Bayesian model selection.
\end{itemize}

\end{document}
